\documentclass[conference]{IEEEtran}
\pagestyle{plain}

% --- Codificación e Idioma ---
\usepackage[utf8]{inputenc}
\usepackage[T1]{fontenc}
\usepackage[spanish,es-nodecimaldot,es-noquoting,es-tabla]{babel}
% Desactivar manejo especial de % de babel spanish
\makeatletter
\def\es@sppercent{}
\makeatother

% --- Microtipografía y Fuentes ---
\usepackage{microtype}

% --- Matemáticas y Unidades ---
\usepackage{amsmath,amssymb,amsfonts}
\usepackage{siunitx}
\sisetup{output-decimal-marker = {.}, per-mode = symbol}
\DeclareSIUnit\rpm{RPM}
\DeclareSIUnit\baud{baudios}
\DeclareSIUnit\baudios{baudios}

% --- Gráficos y Rutas ---
\usepackage{graphicx}
\usepackage[caption=false,font=footnotesize]{subfig}
\graphicspath{{./figuras/}}
\DeclareGraphicsExtensions{.pdf,.jpeg,.png,.PNG,.jpg,.JPG}

% --- Tablas Científicas ---
\usepackage{booktabs}
\usepackage{array}
\usepackage{makecell}

% --- Circuitos y Diagramas Vectoriales TikZ ---
\usepackage{tikz}
\usetikzlibrary{arrows.meta, positioning, fit, calc, shapes.geometric, backgrounds, babel}
\usepackage[siunitx, european]{circuitikz}

% --- Código Fuente Embebido ---
\usepackage{listings}
\usepackage{xcolor}

\definecolor{codeblue}{rgb}{0.13,0.13,0.7}
\definecolor{codegreen}{rgb}{0,0.5,0}
\definecolor{codegray}{rgb}{0.45,0.45,0.45}
\definecolor{codeback}{rgb}{0.97,0.97,0.97}

\makeatletter
\def\lst@makecaption#1#2{%
  \par\addvspace{4pt}%
  \begingroup
    \leftskip=0pt\relax
    \rightskip=0pt\relax
    \fontsize{8pt}{10pt}\selectfont
    \textbf{#1}: #2\par
  \endgroup
}
\makeatother

\lstdefinestyle{embeddedc}{
    backgroundcolor=\color{codeback},
    commentstyle=\color{codegreen}\itshape,
    keywordstyle=\color{codeblue}\bfseries,
    numberstyle=\tiny\color{codegray},
    stringstyle=\color{purple},
    basicstyle=\ttfamily\scriptsize,
    breaklines=true,
    breakatwhitespace=true,
    captionpos=b,
    numbers=left,
    numbersep=5pt,
    xleftmargin=14pt,
    xrightmargin=0pt,
    tabsize=2,
    language=C++
}
\lstset{style=embeddedc}

% --- Hipervínculos y Referencias ---
\usepackage{xurl}
\usepackage{cite}
\usepackage[colorlinks=true, linkcolor=black, citecolor=black, urlcolor=black]{hyperref}

\def\figureautorefname{Figura}
\def\equationautorefname{Ecuación}
\def\tableautorefname{Tabla}
\def\sectionautorefname{Sección}
\def\subsectionautorefname{Subsección}

\IEEEoverridecommandlockouts{}

\tolerance=10000
\hbadness=10000
\emergencystretch=1.5em

\begin{document}

\title{Diseño Mecatrónico, Identificación Experimental y Control Digital Determinista a \qty{1}{\kilo\hertz} para un Péndulo Invertido Rotacional de Furuta con Gemelo Digital 3D}

\author{
    \IEEEauthorblockN{J. Pineda\IEEEauthorrefmark{1}, M. Rodríguez\IEEEauthorrefmark{2} y J. Salazar\IEEEauthorrefmark{3}}
    \IEEEauthorblockA{
        \textit{Escuela de Ingeniería, Facultad de Ciencias Básicas e Ingeniería} \\
        \textit{Universidad de los Llanos} \\
        \textit{Villavicencio, Colombia} \\
        \small \IEEEauthorrefmark{1}\texttt{jspineda.ospina@unillanos.edu.co} \quad (Cód. 161004610) \\
        \small \IEEEauthorrefmark{2}\texttt{marodriguez.payares@unillanos.edu.co} \quad (Cód. 161004834) \\
        \small \IEEEauthorrefmark{3}\texttt{jssalazar.vargas@unillanos.edu.co} \quad (Cód. 161004938)
    }
}

\maketitle

\renewcommand\abstractname{Resumen}
\begin{abstract}
En este trabajo se diseña, caracteriza experimentalmente e implementa en tiempo real el control digital determinista de un péndulo invertido rotacional de Furuta subactuado e inestable en lazo abierto. La plataforma física se manufacturó en PLA mediante diseño para fabricación aditiva (DfAM) sin soportes ($\text{SUP}=0$, $0/23$ interferencias en FreeCAD 1.1) y se instrumentó sobre una tarjeta STM32 Nucleo-F446RE (Cortex-M4F a \qty{180}{\mega\hertz}) con lazo isócrono \texttt{TIM6} a \qty{1.0}{\kilo\hertz}, etapa de potencia L298N modulada por PWM a \qty{20.0}{\kilo\hertz} y dos encoders magnéticos AS5600 de 12~bits leídos por doble bus Soft-$\text{I}^2\text{C}$ ($\sim\qty{450}{\kilo\hertz}$). A partir de $3\,726$ muestras experimentales en banco, se identificaron la dinámica del actuador ($G_v(s)=3.3720/(0.1175s+1)$, $R^2=98.51\%$), la zona muerta asimétrica ($+10.0\% / -21.0\%$~PWM) y los parámetros acoplados pendulares ($\omega_n=\qty{8.8484}{\radian\per\second}$, $\zeta=0.06488$, $K_{\text{acople}}=0.8737$). Con base en ello, se integró el esquema híbrido \texttt{M4 Ultra-Silk}: levantamiento resonante en cuadratura acotado en carrera, lazo dual de equilibrio ($\text{PD}_\alpha + \text{PID}_\theta$), cascada biquad \textit{Dual-Notch} (\qty{11.0}{\hertz} fijo y \qty{6.2}{\hertz} adaptativo), rampa continua compensadora L298N y servo de retorno con reinicio de Clegg y freno activo. En validación física continua (\qty{15.001}{\second}), el prototipo levantó en \qty{0.98}{\second} y sostuvo el equilibrio por \qty{14.02}{\second} con cero caídas ($\mathrm{RMS}(\tilde{\alpha})=\qty{1.99}{\degree}$, $u_{\mathrm{RMS}}=\qty{0.73}{\volt}$), supervisado en tiempo real por un gemelo digital 3D en MATLAB.
\end{abstract}

\renewcommand\IEEEkeywordsname{Palabras clave}
\begin{IEEEkeywords}
Péndulo de Furuta, Control Digital, STM32F446RE, Identificación Experimental, M4 Ultra-Silk, Filtros Dual-Notch, Reinicio de Clegg, DfAM, Gemelo Digital 3D.
\end{IEEEkeywords}

\input{secciones/01_introduccion.tex}
\input{secciones/02_metodologia.tex}
\input{secciones/03_resultados.tex}
\input{secciones/04_discusion.tex}
\input{secciones/05_conclusiones.tex}

\bibliographystyle{IEEEtran}
\bibliography{plantilla}

\end{document}
